\documentclass[12pt]{article}
\usepackage[margin=1in]{geometry}
\usepackage{amsmath}
\usepackage{amssymb}
\usepackage{graphicx}
\usepackage{booktabs}
\usepackage{hyperref}
\usepackage[capitalise]{cleveref}
\usepackage{microtype}
\usepackage{parskip}

\graphicspath{{../refs/images/}}

\newcommand{\figplaceholder}[2][3in]{%
  \begin{center}
  \fbox{\parbox{#1}{\centering\vspace{0.5in}\textit{[Figure: #2]}\vspace{0.5in}}}
  \end{center}
}

\title{Frequency Response, Part~2: Stability Margins and Gain Design\\[0.5em]
\large ME 2801 --- Week 9}
\author{}
\date{}

\begin{document}
\maketitle
\tableofcontents
\newpage

% ===========================================================================
\section{From Bode Plot to Design}
\label{sec:intro}
% ===========================================================================

In Week~8 we built the vocabulary of Bode plots: how to read the frequency
response of a system from its transfer function and how to sketch it using the
building-block rules.  This week we put that vocabulary to work.

The Bode plot of the \emph{open-loop} transfer function carries three pieces of
design information:

\begin{enumerate}
  \item \textbf{Stability.}  The shape of the open-loop Bode plot reveals
    whether the closed-loop system is stable.
  \item \textbf{Percent overshoot.}  A single number read from the Bode
    plot---the \emph{phase margin}---is directly related to the closed-loop
    damping ratio and therefore to percent overshoot.
  \item \textbf{Speed.}  The closed-loop bandwidth is related to a frequency
    we can read from the Bode plot---the gain crossover frequency.
\end{enumerate}

All three of these properties can be adjusted by scaling the open-loop gain.
Multiplying the open-loop transfer function by a constant $K$ shifts the
magnitude curve up or down by $20\log K$~dB without changing the phase curve.
This gain lever is the subject of this chapter.

% ===========================================================================
\section{Stability Margins from the Bode Plot}
\label{sec:margins}
% ===========================================================================

% ---------------------------------------------------------------------------
\subsection{The Nyquist stability criterion}
\label{sec:nyquist}
% ---------------------------------------------------------------------------

The Nyquist stability criterion gives the definitive test for closed-loop
stability from the open-loop frequency response.  The criterion has a full
derivation based on complex analysis that we will not reproduce here.  For the
systems we work with in this course---open-loop stable plants, unity-feedback
loops---the criterion reduces to the following:

\begin{quote}
\textbf{Nyquist stability criterion (restricted form):} A unity-feedback
closed-loop system whose open-loop transfer function $G(s)$ has all poles in
the left half-plane is stable if and only if
$|G(j\omega)| < 1$ (i.e., $< 0$~dB) at every frequency where
$\angle G(j\omega) = -180^\circ$.
\end{quote}

The intuition is straightforward.  A phase of $-180^\circ$ means the feedback
signal has been inverted by the time it returns to the summing junction.  If
the gain around the loop at that frequency also exceeds unity, the inverted
error feeds back positively and the system goes unstable.  The stability
condition says: whenever the loop can invert the signal, it must not amplify
it.

% ---------------------------------------------------------------------------
\subsection{Two key frequencies}
\label{sec:crossover_freqs}
% ---------------------------------------------------------------------------

Two frequencies appear repeatedly in stability-margin analysis.

The \textbf{gain crossover frequency} $\omega_{\Phi_M}$ is the frequency at
which the open-loop magnitude equals 0~dB:
\begin{equation}
  |G(j\omega_{\Phi_M})| = 1
  \label{eq:gc_freq}
\end{equation}

The \textbf{phase crossover frequency} $\omega_{G_M}$ is the frequency at
which the open-loop phase equals $-180^\circ$:
\begin{equation}
  \angle G(j\omega_{G_M}) = -180^\circ
  \label{eq:pc_freq}
\end{equation}

These are the two reference points from which the stability margins are read.

% ---------------------------------------------------------------------------
\subsection{Gain margin and phase margin}
\label{sec:gm_pm}
% ---------------------------------------------------------------------------

\begin{figure}[htbp]
  \figplaceholder[5in]{Bode magnitude (top) and phase (bottom) plots
    annotated with gain margin $G_M$ and phase margin $\Phi_M$.
    On the magnitude plot, $G_M$ is the distance from 0~dB down to the
    magnitude curve at $\omega_{G_M}$.
    On the phase plot, $\Phi_M$ is the distance from $-180^\circ$ up to
    the phase curve at $\omega_{\Phi_M}$. [Nise Fig.~10.37]}
  \caption{Gain and phase margins on the Bode diagram.  Positive margins
    indicate a stable closed-loop system.}
  \label{fig:margins_bode}
\end{figure}

Figure~\ref{fig:margins_bode} shows both margins on a Bode diagram.

\textbf{Phase margin} $\Phi_M$ is how much additional phase lag would be
required to push the system to the edge of instability.  Read from the
\emph{phase} plot at the gain crossover frequency:
\begin{equation}
  \boxed{\Phi_M = 180^\circ + \angle G(j\omega_{\Phi_M})}
  \label{eq:phase_margin}
\end{equation}
Since the phase at the gain crossover of a stable system is between $-90^\circ$
and $-180^\circ$, this formula produces a positive number.

\textbf{Gain margin} $G_M$ is how much additional gain could be added before
instability.  Read from the \emph{magnitude} plot at the phase crossover
frequency:
\begin{equation}
  \boxed{G_M = -20\log|G(j\omega_{G_M})| \quad \text{(dB)}}
  \label{eq:gain_margin}
\end{equation}
A positive gain margin means the magnitude curve is below 0~dB when the phase
reaches $-180^\circ$---the gain must be multiplied by $10^{G_M/20}$ before
instability.

\medskip
\noindent\textbf{Stability rules.}  For an open-loop stable plant:
\begin{itemize}
  \item Positive $G_M$ and positive $\Phi_M$: closed-loop stable.
  \item Negative $G_M$ or negative $\Phi_M$: closed-loop unstable.
\end{itemize}
Both margins must be positive for stability.  A system with either margin equal
to zero is at the boundary of instability (marginally stable).

% ---------------------------------------------------------------------------
\subsection{Worked example: reading margins}
\label{sec:margin_example}
% ---------------------------------------------------------------------------

\textbf{System.}  Consider the unity-feedback system with open-loop transfer
function
\begin{equation}
  G(s) = \frac{200}{(s+2)(s+4)(s+5)}
  \label{eq:ex_loop}
\end{equation}

\begin{figure}[htbp]
  \figplaceholder[5.5in]{Bode magnitude and phase plots for $G(s) = 200/[(s+2)(s+4)(s+5)]$.
    Annotate: gain crossover at $\omega_{\Phi_M} \approx 5.5$~rad/s where
    phase reads $-165^\circ$; phase crossover at $\omega_{G_M} \approx 7$~rad/s
    where magnitude reads $-6$~dB.  Mark the resulting $\Phi_M = 15^\circ$
    and $G_M = 6$~dB. [Nise Fig.~10.36 scaled to $K = 200$]}
  \caption{Bode plots for the margin example.  Phase margin and gain margin
    are read directly from the two crossover frequencies.}
  \label{fig:ex_bode}
\end{figure}

\textbf{Phase margin.}  From Figure~\ref{fig:ex_bode}, the magnitude curve
crosses 0~dB at $\omega_{\Phi_M} \approx 5.5$~rad/s.  At that frequency the
phase reads approximately $-165^\circ$, so:
\[
  \Phi_M = 180^\circ + (-165^\circ) = 15^\circ
\]

\textbf{Gain margin.}  The phase curve reaches $-180^\circ$ at
$\omega_{G_M} \approx 7$~rad/s.  At that frequency the magnitude reads
approximately $-6$~dB, so:
\[
  G_M = -(-6\text{ dB}) = 6\text{ dB}
\]

Both margins are positive: the closed-loop system is stable.  But a phase
margin of $15^\circ$ is small---we expect a lightly damped, oscillatory
transient response.  The next section quantifies this connection.

% ===========================================================================
\section{Connecting Frequency Response to Transient Response}
\label{sec:connect}
% ===========================================================================

% ---------------------------------------------------------------------------
\subsection{Closed-loop frequency response characteristics}
\label{sec:cl_freq}
% ---------------------------------------------------------------------------

The Bode plot is the \emph{open-loop} frequency response.  The \emph{closed-loop}
frequency response is a different curve, and its shape also tells us about the
transient response.

For the standard second-order system with closed-loop transfer function
\begin{equation}
  T(s) = \frac{\omega_n^2}{s^2 + 2\zeta\omega_n s + \omega_n^2}
  \label{eq:cl_tf}
\end{equation}
the closed-loop magnitude $M(\omega) = |T(j\omega)|$ is a function of frequency
that peaks (for $\zeta < 0.707$) at the \textbf{peak frequency}:
\begin{equation}
  \omega_p = \omega_n\sqrt{1 - 2\zeta^2}
  \label{eq:peak_freq}
\end{equation}
The \textbf{peak magnitude} at that frequency is:
\begin{equation}
  M_p = \frac{1}{2\zeta\sqrt{1-\zeta^2}}
  \label{eq:mp}
\end{equation}
For $\zeta \geq 0.707$ the response is monotonically decreasing with no peak
above unity.

\begin{figure}[htbp]
  \figplaceholder[4.5in]{Closed-loop log-magnitude frequency response for the
    standard second-order system.  Mark $M_p$ and $\omega_p$; mark the
    $-3$~dB bandwidth frequency $\omega_{BW}$. [Nise Fig.~10.39]}
  \caption{Closed-loop magnitude response showing peak $M_p$, peak frequency
    $\omega_p$, and bandwidth $\omega_{BW}$.}
  \label{fig:cl_mag}
\end{figure}

The \textbf{bandwidth} $\omega_{BW}$ is the $-3$~dB frequency of the
closed-loop magnitude response (Figure~\ref{fig:cl_mag}):
\begin{equation}
  \omega_{BW} = \omega_n\sqrt{(1-2\zeta^2)+\sqrt{4\zeta^4-4\zeta^2+2}}
  \label{eq:bw}
\end{equation}
Bandwidth measures the speed of the closed-loop frequency response: larger
bandwidth means faster time-domain response.  Using the formulas for $\omega_n$
in terms of settling time and peak time, Eq.~\eqref{eq:bw} can be rewritten as:
\begin{align}
  \omega_{BW} &= \frac{4}{T_s\zeta}
    \sqrt{(1-2\zeta^2)+\sqrt{4\zeta^4-4\zeta^2+2}}
  \label{eq:bw_ts} \\[4pt]
  \omega_{BW} &= \frac{\pi}{T_p\sqrt{1-\zeta^2}}
    \sqrt{(1-2\zeta^2)+\sqrt{4\zeta^4-4\zeta^2+2}}
  \label{eq:bw_tp}
\end{align}
Given a speed specification (settling time or peak time) and a damping ratio,
these formulas give the required closed-loop bandwidth.

% ---------------------------------------------------------------------------
\subsection{Phase margin and damping ratio}
\label{sec:pm_zeta}
% ---------------------------------------------------------------------------

We can compute the Bode plot of the \emph{open-loop} transfer function easily
by hand or with MATLAB\@.  The closed-loop frequency response requires more work.
Phase margin provides the bridge: it connects the open-loop Bode plot directly
to the closed-loop damping ratio.

For the second-order open-loop transfer function
\begin{equation}
  G(s) = \frac{\omega_n^2}{s(s+2\zeta\omega_n)}
  \label{eq:second_order_ol}
\end{equation}
the phase margin and damping ratio are related by:
\begin{equation}
  \boxed{
    \Phi_M = \tan^{-1}\!\left(\frac{2\zeta}{\sqrt{-2\zeta^2+\sqrt{1+4\zeta^4}}}\right)
  }
  \label{eq:pm_zeta}
\end{equation}

\begin{figure}[htbp]
  \figplaceholder[3.5in]{Phase margin $\Phi_M$ (degrees) vs.\ damping ratio
    $\zeta$ from Eq.~\eqref{eq:pm_zeta}.  The curve is monotonically
    increasing from $0^\circ$ at $\zeta=0$ through $65.5^\circ$ at
    $\zeta=0.707$ and approaching $90^\circ$ for large $\zeta$.
    [Nise Fig.~10.48]}
  \caption{Phase margin vs.\ damping ratio for the standard second-order
    open-loop transfer function.}
  \label{fig:pm_zeta}
\end{figure}

Figure~\ref{fig:pm_zeta} shows that the relationship is monotonic: more damping
always produces a larger phase margin.  This gives us the complete design chain:

\begin{equation*}
  \text{\%OS specification}
  \;\xrightarrow{\text{Eq.~(4.39)}}\;
  \zeta
  \;\xrightarrow{\text{Eq.~\eqref{eq:pm_zeta}}}\;
  \Phi_M \text{ (required)}
  \;\xrightarrow{\text{Bode plot}}\;
  K \text{ (designed)}
\end{equation*}

Specify percent overshoot, convert to damping ratio, convert to required phase
margin, and then use the Bode plot to find the gain that achieves it.

\medskip
\noindent\textbf{Note on generality.}  Equation~\eqref{eq:pm_zeta} is exact for the
pure second-order open-loop $G(s)$.  For higher-order systems, the formula is
an approximation---it becomes less accurate when additional poles and zeros
significantly alter the phase near the gain crossover frequency.  We use it as
a design starting point and verify with simulation.

% ---------------------------------------------------------------------------
\subsection{Bandwidth from the open-loop Bode plot}
\label{sec:bw_ol}
% ---------------------------------------------------------------------------

We showed that the closed-loop bandwidth $\omega_{BW}$ is related to the
transient response speed via Eqs.~\eqref{eq:bw_ts} and~\eqref{eq:bw_tp}.  We
can also estimate $\omega_{BW}$ directly from the open-loop Bode plot without
computing the closed-loop frequency response.

The key observation (from the Nichols chart analysis) is:

\begin{quote}
The closed-loop bandwidth $\omega_{BW}$ approximately equals the frequency
at which the open-loop magnitude is between $-6$ and $-7.5$~dB, provided the
open-loop phase at that frequency is between $-135^\circ$ and $-225^\circ$.
\end{quote}

In practice: scan the open-loop magnitude plot for the frequency where the
curve is in the $-6$ to $-7.5$~dB band while the phase is in the stated
range.  That frequency is a good estimate of the closed-loop bandwidth.

Combined with the phase margin--damping ratio relationship, this gives a
complete picture of the closed-loop transient response from the open-loop
Bode plot alone:
\begin{itemize}
  \item Phase margin at gain crossover $\rightarrow$ $\zeta$ $\rightarrow$ \%OS
  \item Open-loop magnitude at $-6$ to $-7.5$~dB $\rightarrow$ $\omega_{BW}$ $\rightarrow$ $T_s$ or $T_p$
\end{itemize}

% ===========================================================================
\section{Gain Adjustment Design}
\label{sec:gain_design}
% ===========================================================================

% ---------------------------------------------------------------------------
\subsection{How gain affects the Bode plot}
\label{sec:gain_effect}
% ---------------------------------------------------------------------------

When we scale the open-loop transfer function by a constant $K$:
\begin{itemize}
  \item The magnitude curve shifts \emph{up} by $20\log K$~dB (for $K > 1$)
    or \emph{down} (for $K < 1$).
  \item The phase curve is \emph{unchanged}.
\end{itemize}

Shifting the magnitude curve up moves the gain crossover frequency to a
higher value, which generally:
\begin{itemize}
  \item Increases the closed-loop bandwidth (faster response).
  \item Decreases the phase margin (more overshoot, less damping).
\end{itemize}
Shifting the magnitude curve down has the opposite effects.  By choosing $K$
we trade off speed against damping.

% ---------------------------------------------------------------------------
\subsection{Design procedure}
\label{sec:procedure}
% ---------------------------------------------------------------------------

To design for a percent overshoot specification:

\begin{enumerate}
  \item \textbf{Draw.}  Plot the Bode magnitude and phase diagrams for a
    convenient starting gain.
  \item \textbf{Compute.}  From the percent overshoot requirement, find $\zeta$
    using
    \[
      \%OS = 100\,e^{-\pi\zeta/\sqrt{1-\zeta^2}}
    \]
    Then find the required phase margin from Eq.~\eqref{eq:pm_zeta}.
  \item \textbf{Find.}  On the phase plot, locate the frequency $\omega^*$
    where
    \[
      \angle G(j\omega^*) = \Phi_M - 180^\circ
    \]
    This is the frequency that must become the new gain crossover.
  \item \textbf{Adjust.}  Read the current magnitude at $\omega^*$.  Change
    $K$ to move the magnitude curve so that it crosses 0~dB at $\omega^*$.
    The required adjustment in dB is:
    \[
      \Delta K_{\mathrm{dB}} = -20\log|G(j\omega^*)|_{\mathrm{current}}
    \]
    and the final gain is:
    \[
      K = K_{\mathrm{start}} \times 10^{\Delta K_{\mathrm{dB}}/20}
    \]
\end{enumerate}

% ---------------------------------------------------------------------------
\subsection{Worked example: position control system}
\label{sec:design_example}
% ---------------------------------------------------------------------------

\textbf{System.}  A position control system (preamplifier gain $K$, power
amplifier, motor, and output integrator) has the open-loop transfer function:
\begin{equation}
  G(s) = \frac{100K}{s(s+36)(s+100)}
  \label{eq:ex2_loop}
\end{equation}
\textbf{Specification.}  Design for 9.5\% overshoot.

\begin{figure}[htbp]
  \figplaceholder[5.5in]{Bode magnitude and phase for $G(s) = 100K/[s(s+36)(s+100)]$.
    Show the initial curve ($K = 3.6$, magnitude starting at 0~dB near
    $\omega = 0.1$~rad/s) and the gain-adjusted curve ($K = 583.9$, shifted
    up by 44.2~dB).  Mark $\omega^* = 14.8$~rad/s on the phase plot where
    phase $= -120.8^\circ$ and on the magnitude plot where the adjusted curve
    crosses 0~dB. [Nise Fig.~11.3]}
  \caption{Bode plots for the gain-adjustment design.  Initial gain (solid)
    and gain-adjusted result (dashed).}
  \label{fig:ex_design}
\end{figure}

\textbf{Step 1: Initial Bode plot.}  Choose $K = 3.6$ as a convenient starting
point that places the initial 0~dB crossing near $\omega = 0.1$~rad/s
(Figure~\ref{fig:ex_design}).

\textbf{Step 2: Required phase margin.}

A 9.5\% overshoot gives $\zeta = 0.6$.  Substituting into
Eq.~\eqref{eq:pm_zeta}:
\[
  \Phi_M = \tan^{-1}\!\left(\frac{2(0.6)}{\sqrt{-2(0.6)^2+\sqrt{1+4(0.6)^4}}}\right) = 59.2^\circ
\]

\textbf{Step 3: Target frequency.}

The new gain crossover must be at the frequency where the phase equals
$59.2^\circ - 180^\circ = -120.8^\circ$.  From the phase plot,
$\omega^* = 14.8$~rad/s.

\textbf{Step 4: Adjust $K$.}

At $\omega^* = 14.8$~rad/s with $K = 3.6$, the magnitude is $-44.2$~dB.
Raising by 44.2~dB:
\[
  K = 3.6 \times 10^{44.2/20} = 3.6 \times 162.2 = 583.9
\]

\textbf{Result.}  The designed open-loop transfer function is:
\begin{equation}
  G(s) = \frac{58{,}390}{s(s+36)(s+100)}
  \label{eq:ex2_result}
\end{equation}
A closed-loop simulation confirms 9.48\% overshoot and $T_p = 0.18$~s.

\medskip
\noindent\textbf{Interpretation.}  The gain adjustment moved the gain crossover
from a low frequency to $\omega^* = 14.8$~rad/s.  At that frequency the phase
was already at $-120.8^\circ$, giving the exact phase margin needed to meet
the overshoot specification.

% ===========================================================================
\section{MATLAB for Stability Margins and Gain Design}
\label{sec:matlab}
% ===========================================================================

The Control System Toolbox provides direct support for all of the operations in
this chapter.

\texttt{margin(sys)} computes and plots the gain and phase margins, annotating
the Bode diagram with both crossover frequencies.  The values can also be
captured numerically:
\begin{verbatim}
[Gm, Pm, Wcg, Wcp] = margin(sys);
\end{verbatim}
where \texttt{Gm} is the gain margin (as a linear ratio, not dB), \texttt{Pm}
is the phase margin in degrees, \texttt{Wcg} is the phase crossover frequency
$\omega_{G_M}$, and \texttt{Wcp} is the gain crossover frequency
$\omega_{\Phi_M}$.

The companion script (in \texttt{code/freq\_response\_p2\_companion.m}) covers:
\begin{itemize}
  \item Computing margins with \texttt{margin()} and reading them from the plot.
  \item Plotting the phase margin--damping ratio relationship from
    Eq.~\eqref{eq:pm_zeta}.
  \item An automated implementation of the four-step gain-adjustment procedure.
  \item Verification of the design via closed-loop step response.
  \item Application to the USV surge model.
\end{itemize}

% ===========================================================================
\section{Summary}
\label{sec:summary}
% ===========================================================================

Table~\ref{tab:summary} collects the key formulas from this chapter.

\begin{table}[htbp]
  \centering
  \caption{Key relationships for frequency-domain stability margins and gain-adjustment design}
  \label{tab:summary}
  \renewcommand{\arraystretch}{1.6}
  \begin{tabular}{@{}p{3.8cm}p{6.5cm}p{3.2cm}@{}}
    \toprule
    Quantity & Formula & Notes \\
    \midrule
    Phase margin
      & $\Phi_M = 180^\circ + \angle G(j\omega_{\Phi_M})$
      & read at 0-dB crossing \\
    Gain margin
      & $G_M = -20\log|G(j\omega_{G_M})|$ (dB)
      & read at $-180^\circ$ crossing \\
    Phase margin from $\zeta$
      & $\Phi_M = \tan^{-1}\!\!\left(\dfrac{2\zeta}{\sqrt{-2\zeta^2+\sqrt{1+4\zeta^4}}}\right)$
      & exact for 2nd-order $G$ \\
    Closed-loop bandwidth
      & $\omega_{BW} = \omega_n\sqrt{(1-2\zeta^2)+\sqrt{4\zeta^4-4\zeta^2+2}}$
      & $-3$\,dB of $|T(j\omega)|$ \\
    $\omega_{BW}$ from $T_s$
      & $\omega_{BW} = \dfrac{4}{T_s\zeta}\sqrt{(1-2\zeta^2)+\sqrt{4\zeta^4-4\zeta^2+2}}$
      & \\
    $\omega_{BW}$ from $T_p$
      & $\omega_{BW} = \dfrac{\pi}{T_p\sqrt{1-\zeta^2}}\sqrt{(1-2\zeta^2)+\sqrt{4\zeta^4-4\zeta^2+2}}$
      & \\
    \bottomrule
  \end{tabular}
\end{table}

The central thread of this chapter: the open-loop Bode plot carries design
information about the closed-loop transient response.  The gain crossover
frequency sets the bandwidth.  The phase at the gain crossover---the phase
margin---determines the damping ratio.  A simple gain adjustment shifts the
magnitude curve, moves the crossover frequency, and thereby sets the phase
margin.  Starting from a percent overshoot specification and working through
the chain $\text{\%OS} \to \zeta \to \Phi_M \to \omega^* \to K$, we can
design the proportional gain directly from the Bode plot.

\end{document}
